\documentclass[mlmain,onecolumn]{jmlr}

\usepackage{booktabs}
\usepackage{multirow}
\usepackage{amsmath}

\newcommand{\auc}{\mathrm{AUC}}

%%%% DON'T CHANGE %%%%%%%%%
\jmlrvolume{}
\firstpageno{1}
\editors{}
\jmlryear{2026}
\jmlrworkshop{Symmetry and Geometry in Neural Representations}
%%%%%%%%%%%%%%%%%%%%%%%%%%%

\title[Cortical Decoders Age Slower]{Cortical Decoders Trained on Other
Individuals Age Slower Than Your Own}

\author{\Name{Anonymous Author(s)}}

\begin{document}
\maketitle

\begin{abstract}
With adapting brain--computer interface (BCI) technology, neural decoding algorithms are one of the
most essential parts of being able to understand people's brains. Decoders must work for months, and
recalibrating them per subject has high subject costs. A decoder fit to a singular animal's cortex
drifts, losing accuracy over days, and usually the remedy is recalibrating on fresh data from that
same subject. We report a different kind of robustness: a
decoder built from one other individual, or pooled from a set of other individuals, ages more slowly
than the subject's own. On the open BraiDyn-BC cohort (25 mice, 44 Allen-atlas cortical parcels, a
$\sim$15-day operant protocol) we decode four task events and measure, per mouse, the $\auc$ lost
over two weeks by its own early-week decoder against one pooled over the other 24 mice and one built
from a single other mouse. The personal decoder ages more by $+0.017$ $\auc$
($p = 5 \times 10^{-5}$, 73 of 100 mouse--event pairs). This also is not a data-volume effect: at a matched training-event count a
decoder from a single other mouse still ages less. Hence we conclude this is rather an effect of not
overfitting an individual's drift-prone idiosyncrasies. It also survives a 1-D CNN and a GRU. The
effect holds past this cohort and recording modality. On an independent two-photon dataset from a
different laboratory -- Allen Visual Behavior, 23 mice, single-plane VISp populations decoding a
stimulus change -- the same session-held-out protocol gives $+0.023$ $\auc$
($p = 0.003$, 74\% of mice), and the count-matched control reproduces there as well
(self $-$ one other mouse $= +0.011$ $\auc$). The
early-week accuracy cost of decoding across individuals is repaid as late-week stability, which
suggests population priors as an alternative to per-subject recalibration. A second offering is the
session-held-out protocol under which all of this is measured: because only the within-subject arm
can share a recording session with its test data, a trial-level split inflates that arm alone, and
holding out whole sessions is what allows this form of robustness to be reported reliably. Code and
a streamed-feature pipeline are released.
\end{abstract}

\section{Introduction}
Neural decoding is one of the most essential ways to understand neural interfaces and create BCI
technology, and a decoder that is deployed is a decoder that has to keep working: over months, on
the same individual, without fresh labelled data arriving to prop it up. Representational drift
stands directly against that. The reorganization of an animal's neural code over days with fixed
behavior is now documented across sensory, motor, and association cortex and in the hippocampus
\citep{rule2020,driscoll2017,schoonover2021}. This has a stark impact on neural decoding: a decoder
trained on an animal today degrades on that same animal weeks later. The dominant response is
recalibration, periodically refitting the decoder on fresh within-subject data
\citep{degenhart2020,jarosiewicz2015}, which requires ongoing labelled data from the individual and
treats each subject in isolation.

The code's cross-individual structure offers an alternate route to temporal robustness, and the
field already relies on it. Cortex-wide behavioral representations are considerably shared among
individuals \citep{musall2019,macdowell2020,macdowell2024}; latent dynamics are preserved across
animals well enough that a decoder fit to one transfers to another \citep{safaie2023}; and pooled
decoders already read out behavior from unseen individuals \citep{wicat2026,neds2025,peterson2021}.
Methods such as Meta-AlignNN \citep{metaalign2025} are built on that premise. While people currently
know that conservation across individuals is established, and that pooling is known to work at a
fixed time, every one of these results is measured at a fixed time. We offer pooling as a solution
to a new problem: the fact that it resists drift over time.

Our core discovery is that a decoder that never saw the subject cannot fit to that subject's
drift-prone idiosyncrasies, so it should age more slowly, and the way we measure that difference is
as an asymmetry between the two arms. This prediction, to our knowledge, has not been examined.

We test it on BraiDyn-BC \citep{kondo2025}, a widefield calcium-imaging resource whose $\sim$15-day
operant protocol and shared 44-region Allen parcellation make within-subject and cross-subject
decoders directly comparable across time. Our contributions:
\begin{itemize}\itemsep2pt
\item \textbf{Pooling resists drift.} Across a two-week period the population-pooled decoder loses
less accuracy than the individual's own, with a paired asymmetry of $+0.017$ $\auc$
($p = 5 \times 10^{-5}$, 73 of 100 pairs; Section~\ref{sec:asym}). The size of this effect varies
among the four events, but those are chained within a trial and have highly differing counts, so we
treat the pooled estimate as the result and leave the per-event ranking uninterpreted
(Section~\ref{sec:events}).
\item \textbf{It is not a data-volume artifact.} At a matched training-event count a decoder trained
on one other mouse ages less than one trained on the subject's own, by $+0.010$ $\auc$ on average.
The gain therefore comes from training on individuals other than the subject, not from the larger
training set pooling usually brings, and training on many other mice rather than one lowers aging
further on all four events (Sections~\ref{sec:matched}--\ref{sec:scaling}).
\item \textbf{It replicates on an independent dataset.} On the Allen Visual Behavior two-photon
cohort the same protocol gives $+0.023$ $\auc$ ($p = 0.003$, 74\% of 23 mice), with the
count-matched ordering also reproduced (Section~\ref{sec:allen}).
\item \textbf{Session-level validation is load-bearing.} As trials within a recording session are
dependent, a trial-level split inflates the personal decoder but not the pooled one, one-sidedly
inflating the asymmetry. We report that inflation for each event; on lever pull it accounts for the
entire apparent effect. For longitudinal decoding comparisons we recommend session-held-out
validation as a default (Section~\ref{sec:protocol}).
\end{itemize}

\section{Related work}
Drift is characterized within individual subjects
\citep{rule2020,driscoll2017,schoonover2021,rule2019}, and most approaches to mitigating it also
operate within the subject: manifold stabilization \citep{degenhart2020}, self-calibrating
recalibration \citep{jarosiewicz2015,wilson2025}, and alignment to the subject's own reference day
\citep{farshchian2019,karpowicz2025}. Each uses later data from the same subject. Cross-subject
structure separately supports readout \citep{melbaum2022,schneider2023,wicat2026} and reduces
overfitting when pooled \citep{peterson2021,neds2025}, and latent dynamics are preserved across
animals well enough that a decoder transfers between them \citep{safaie2023}. That conservation is
why pooling is possible here; it does not imply pooling will \emph{resist drift}, since a code can be
shared across individuals and still move within each at the same rate. The closest work,
Meta-AlignNN \citep{metaalign2025}, proposes a method for cross-subject stability, whereas we measure
whether a difference in aging rate exists at all, with training-set size held fixed.
Appendix~\ref{app:related} gives the full discussion.

\section{Data and methods}
\subsection{Datasets and preprocessing}
BraiDyn-BC \citep{kondo2025} (DANDI:001425, CC-BY) provides, for each mouse and
session, hemodynamically corrected dorsal-cortex $\Delta F/F$ traces parcellated into 44 Allen Common
Coordinate Framework regions at $\sim$30\,Hz, together with time-aligned behavioral channels. The
dataset includes 25 mice, each of which completed $\sim$15 daily operant sessions over $\sim$3 weeks.
We stream the 44-region traces and behavioral channels without downloading the raw movies.

The two cohorts used in this paper were built for different purposes, and the ways they differ are
what make the comparison informative rather than redundant. BraiDyn-BC is densely longitudinal, with
roughly fifteen sessions per mouse over three weeks, and it supplies a shared coordinate system: the
44 Allen parcels mean that region $k$ in one mouse is the same region in every other, so a pooled
decoder can be fit directly in a common feature space. The Allen Visual Behavior cohort has neither
property. Its longitudinal subset is sparser, and because it is cellular two-photon imaging there is
no correspondence between animals at all -- one mouse's cell 17 is unrelated to another's -- so a
common feature space has to be constructed rather than assumed (Section~\ref{sec:allencohort}). The
recording modality, the task, and the decoded contrast all differ as well. A result that survives
both is therefore not a property of one atlas, one instrument, or one behavior.

\subsection{Features and decoders}
We detect the onset of four events (lever pull, lick, reward, and
tone) as rising edges, with a 1\,s refractory period. For each onset we construct a 44-dimensional
evoked feature by subtracting the mean pre-onset baseline ($-0.5$--$0$\,s) from the mean post-onset
$\Delta F/F$ response ($0$--$1$\,s) in each parcel. Negative samples are drawn from matched-count
windows more than 2\,s from any onset \emph{of the same event}; they are not screened against the
other three events, and Section~\ref{sec:events} reports how often they contain one. Every
classification is a binary decode of one event against these matched negatives, scored by area under
the ROC curve ($\auc$). The four decoders are therefore four independent binary detectors, not a
four-way classifier: nothing in this paper asks a model to distinguish one event from another.

The primary decoder is a standardized $\ell_2$-regularized logistic regression on
the 44-dimensional evoked feature. This is a linear readout of the parcel-space code: it recovers the
component of behavior that is linearly present in the population average, and it is the natural
comparison to prior cross-subject decoders, which are overwhelmingly linear
\citep{melbaum2022,peterson2021}. Every condition below uses this same readout, so the only variable
is which data the decoder was trained on.\footnote{The result is not a property of the linear
readout. Repeating the analysis with a 1-D convolutional network over time and a GRU over the frame
sequence (full evoked window, $-0.5$ to $+1.0$\,s, $45 \times 44$), under the same session-held-out
protocol, the asymmetry survives for both events that carry it: lick gives $+0.020$, $+0.012$ and
$+0.010$ and reward $+0.029$, $+0.015$ and $+0.012$ under the linear, CNN and GRU readouts, all
significant. Under a trial-level split the nonlinear estimates appear \emph{larger} instead, which is
the inflation of Section~\ref{sec:protocol}: a higher-capacity model exploits it more effectively
than a linear one.}

\subsection{Temporal blocks, aging, and the session-held-out protocol}
\label{sec:blocks}
We group each mouse's sessions into an early block (days 1--5)
and a late block (days 11--15) and report $\auc$ under four conditions. The personal decoder is
estimated by leaving out one whole early \emph{session} at a time: for each early session $s$ we
train on the mouse's remaining early sessions and score the resulting decoder both on $s$
(within-mouse same-block, WS) and on the late block (within-mouse across-block, WX), then average
over $s$. Because WS and WX come from the same fitted decoders, they differ only in which data they
are tested on. The pooled decoder is trained once on the other mice's early blocks and scored on the
same held-out sessions $s$ (LS) and on the late block (LX).

Holding out sessions rather than trials is essential and not a detail. Trials within one recording
share illumination and hemodynamic state, slow arousal drift, and temporal autocorrelation, so a
trial-level split lets the personal decoder train and test inside the same session. The pooled
decoder is trained on other animals entirely and cannot benefit this way, so the inflation is
one-sided and lands directly on the quantity of interest; Section~\ref{sec:protocol} reports what it
costs per event.

Aging is the accuracy lost when testing moves from the held-out mouse's early data to its late block
while training stays fixed: personal aging is $\mathrm{WS}-\mathrm{WX}$ and pooled aging is
$\mathrm{LS}-\mathrm{LX}$. The per-mouse aging asymmetry is
$(\mathrm{WS}-\mathrm{WX})-(\mathrm{LS}-\mathrm{LX})$; a positive value means the personal decoder
ages more. We assess significance with a paired sign-flip permutation test across mice
(20{,}000 permutations) and compute 95\% confidence intervals with a cluster bootstrap over mice.

\subsection{The external validation cohort}
\label{sec:allencohort}
Section~\ref{sec:allen} repeats the analysis on the Allen Visual Behavior two-photon dataset
\citep{allenvb2023}: containers spanning at least 14 days whose experiments are all in VISp within a
single 50\,\textmu m depth bin (the largest such band, 150--199\,\textmu m, holds 25 mice and 184
experiments), leaving 166 experiments and 23 mice after per-experiment quality control. Because
different animals have no neuron correspondence and the longitudinal subset is single-plane,
averaging over the field of view would give a one-dimensional feature; the representation that is
comparable across animals is therefore a permutation-invariant summary of the population response,
giving a fixed 27-dimensional feature regardless of how many cells were segmented. The decoded
contrast is \texttt{is\_change}. Appendix~\ref{app:allen} gives the full selection rule and feature
definition.

\section{Results}

\subsection{The effect: a pooled decoder ages more slowly, by two different estimators}
\label{sec:asym}
\label{sec:decay}
For each mouse we keep the test data fixed (first the mouse's held-out early sessions, then its late
block) and change only the source of the training data: either the mouse itself or the other 24 mice.
Across all 100 mouse--event pairs the asymmetry is $+0.017$ $\auc$ ($p = 5 \times 10^{-5}$), with the
personal decoder aging more in 73\% of pairs (Table~\ref{tab:asym}).

The effect differs among the four events: it is largest for reward ($+0.029$, $p = 0.0005$) and lick
($+0.020$, $p = 0.002$), marginal for tone ($+0.015$, $p = 0.08$), and absent for lever pull
($+0.003$, $p = 0.37$). While we report this variation, we do not interpret it, as the four events
all fall within roughly the same two seconds of a trial. Accordingly, the pooled estimate over all
100 pairs is the claim we defend, and the per-event column should be read as data rather than as
four separate findings.

\begin{table}[t]\centering
\caption{Personal vs.\ pooled aging across days ($\auc$ lost from early to late test, holding
training fixed), per event, $n = 25$ mice, with whole early sessions held out. Asymmetry $=$
personal $-$ pooled aging; a positive value means the mouse's own decoder ages more. $p$: paired
sign-flip permutation. The rightmost column gives the same asymmetry under a trial-level split, the
protocol this analysis replaces (Section~\ref{sec:protocol}).}
\label{tab:asym}
\setlength{\tabcolsep}{4pt}
\begin{tabular}{lccccccc}
\toprule
Event & Personal & Pooled & Asymmetry (95\% CI) & $p$ & own ages more & trial-split \\
\midrule
Lever-pull & $-0.022$ & $-0.025$ & $+0.003$ $[-0.004, 0.011]$ & 0.37 & 64\% & $+0.014$ \\
Lick       & $+0.012$ & $-0.009$ & $+0.020$ $[0.008, 0.032]$ & 0.0023 & 80\% & $+0.028$ \\
Reward     & $+0.017$ & $-0.012$ & $+0.029$ $[0.014, 0.046]$ & 0.0005 & 80\% & $+0.032$ \\
Tone       & $+0.016$ & $+0.001$ & $+0.015$ $[0.001, 0.032]$ & 0.077 & 68\% & $+0.022$ \\
\midrule
\textbf{Pooled} & & & $\mathbf{+0.017}$ $[0.010, 0.024]$ & $\mathbf{5\times10^{-5}}$ & \textbf{73\%} & $+0.024$ \\
\bottomrule
\end{tabular}
\end{table}

The effect also appears as a continuous trend across all fifteen days rather than only as a
two-block contrast. The personal decoder starts ahead by about $0.035$ $\auc$ on day 1, and the gap
narrows by $0.0014$ $\auc$ per day ($p = 5 \times 10^{-5}$; narrowing in 65\% of pairs), so that by
day 15 roughly 60\% of the early-week advantage is gone. Averaged over the four events the pooled
decoder \emph{gains} accuracy over days ($+0.0015$ $\auc$/day) as the animal becomes more practised,
while the personal decoder stays flat ($+0.0000$ $\auc$/day): fit to idiosyncrasies, it is unable to
ride that wave, and the lift it fails to take is what opens the gap. Appendix~\ref{app:decay} gives
the per-event regressions and the day-by-day figure.

\subsection{Ruling out the alternative explanations}
\label{sec:events}
\label{sec:matched}
\label{sec:scaling}
The four decoded events come from a single cued lever-pull trial structure, so they are chained in
time rather than independent. Every reward onset has a lever-pull onset a median of $0.083$\,s away
and a lick onset $0.117$\,s away, and $100\%$ and $98\%$ respectively fall inside the window from
which a feature vector is built. Across all four events, the fraction of positive windows containing
at least one other event is $100\%$ (reward), $84\%$ (lick), $84\%$ (tone) and $68\%$ (lever pull).
Consequently, the reward decoder is never reading reward in isolation: its input window always
contains licking and lever movement as well. The negatives are not event-free either: screened only against
the event's own onsets, they contain another event $33$--$57\%$ of the time.

The events also differ heavily in sample size, from a median of $56$ reward onsets per session
against $107$ (tone), $530$ (lick) and $652$ (lever pull). The per-event asymmetry runs roughly
opposite to that count; a decoder fit to fewer events has more idiosyncrasy available to overfit, and
hence more for pooling to avoid. With only four events, however, we cannot separate that explanation from a
behavioral one, and we therefore do not attempt to. This is why the pooled estimate over all 100 pairs is what we defend, and
why later sections that select reward and lick are consistency checks on the same events rather than
evidence from independent behaviors.

A pooled decoder contains $\sim$24$\times$ more training events than a personal decoder, so its
stability could reflect data volume rather than pooling. We test this with a count-matched control
(Fig.~\ref{fig:mech}, left). For each held-out early session $s$ of each mouse we set the
training-set size to $n$, the number of events in that mouse's remaining early sessions, and compare
three decoders: \emph{self}, trained on those $n$ events; \emph{one}, trained on $n$ events from a
single other mouse; and \emph{many}, trained on $n$ events drawn across the other 24 mice. All three
are scored on the same held-out session $s$ and the same late block, so they differ only in whose
data they were fit on.

For the events where there was an asymmetry, the count-matched ordering reproduces it. The self
decoder ages more than a decoder trained on one other mouse by $+0.010$ $\auc$ on average, and by
$+0.007$ for lick, $+0.017$ for reward and $+0.006$ for tone individually. As both decoders see the
same number of training events, this is not a data-volume artifact.

The diversity ordering is the most robust result in this analysis: \emph{many} ages less than
\emph{one} for all four events, including lever pull ($-0.026$ vs.\ $-0.020$). Training on a diverse
set of individuals rather than a single other individual therefore provides a smaller, additional
benefit that does not depend on which event is decoded.

We next vary the number of training mice, $N \in \{1, 2, 4, 8, 16, 24\}$ (Fig.~\ref{fig:mech},
right). As $N$ increases, pooled aging becomes slightly more favorable for lick ($-0.001$ to
$-0.010$), reward ($-0.004$ to $-0.010$), and tone ($+0.003$ to $-0.010$); lever pull saturates with
a single training mouse and remains near $-0.025$. The scaling effect is around $0.01$ $\auc$,
comparable to the self-versus-one-other gap of Section~\ref{sec:matched}. Most of the resistance to
drift is already present when training on a single other individual, and additional individuals
refine it further. This analysis never lets the decoder see the held-out mouse, so it is unaffected
by the protocol issue below.

\begin{figure}[t]\centering
\includegraphics[width=0.98\linewidth]{figs/fig_mechanism.pdf}
\caption{\textbf{Left:} count-matched control. With equal numbers of training events, decoder aging
(AUC lost from early to late; lower values indicate greater robustness) is compared for the subject's
own data (self), one other mouse, and many other mice. \textbf{Right:} pooled aging as a function of
the number of training mice $N$. Greater diversity provides a small additional gain, with immediate
saturation for lever pull.}
\label{fig:mech}\end{figure}

\subsection{Measuring it correctly: why the validation protocol matters}
\label{sec:protocol}
Every number reported above holds out whole recording sessions, and it is worth showing what that
choice is worth. Repeating the identical analysis on the identical features with a trial-level split -- shuffling trials pooled across the early sessions, so that a
decoder may train on one part of a session and be tested on another -- inflates the personal
decoder's within-block accuracy by $+0.013$ (lever pull), $+0.010$ (lick), $+0.009$ (reward) and
$+0.000$ $\auc$ (tone). The pooled decoder is unaffected, because it is trained on other animals and
has no session in common with the test data.

The inflation therefore lands on one side of a difference between two arms (rightmost column of
Table~\ref{tab:asym}). For lever pull it is the entire effect: a significant-looking $+0.014$
($p = 0.001$) becomes $+0.003$ ($p = 0.37$). For reward, which has the largest true effect, it
changes little. The same reasoning applies to the count-matched control, where the natural
implementation is to score the self decoder on the block it was fit on -- not a cross-validated
quantity at all, and enough to reverse the ordering the control exists to establish.

The mechanism behind this, and why it generalizes beyond our dataset, is set out in
Section~\ref{sec:blocks}.

\subsection{The effect replicates on an independent dataset}
\label{sec:allen}
To further test whether the asymmetry is a property of that one cohort, modality and task, we
repeat the analysis on the Allen Visual Behavior two-photon dataset \citep{allenvb2023}: a different
laboratory, cellular two-photon rather than mesoscale widefield imaging, a visual change-detection
task rather than a cued lever pull, and a 27-dimensional permutation-invariant population summary
rather than a 44-parcel atlas feature. The decoded contrast is a stimulus change, so it is much less
coupled to reward and licking than anything decoded on BraiDyn-BC.

The asymmetry reproduces at a comparable magnitude (Table~\ref{tab:allen}). Across 23 mice the
personal decoder ages by $+0.024$ $\auc$ and the pooled decoder by $+0.001$, an asymmetry of
$+0.023$ $[0.010, 0.037]$, $p = 0.003$, with the mouse's own decoder ageing more in 74\% of animals.
Early-block accuracy is $0.672$ for the personal decoder and $0.638$ for the pooled one, so the
familiar ordering holds at the start; by the late block the personal decoder leads by only $0.012$,
about a third of its initial advantage.

The mechanism identified is not identical in both datasets; on BraiDyn-BC the gap narrows
mostly because the pooled decoder \emph{gains} over days as the animal becomes more practised
(Section~\ref{sec:decay}), while the personal decoder stays flat. Here the pooled decoder is flat
($+0.001$) and the entire asymmetry comes from the personal decoder losing accuracy. The shared
finding within both, hence, is that the personal decoder ages faster.

The count-matched control reproduces as well. At equal training-event count the self decoder ages
by $+0.024$, a decoder from one other mouse by $+0.013$ and one drawn across the other mice by
$+0.003$. The self-minus-one gap is $+0.011$ $[0.001, 0.022]$, close to the $+0.010$ measured on
BraiDyn-BC, though on 23 animals it is marginal by the sign-flip test ($p = 0.054$); self-minus-many
is $+0.021$ $[0.010, 0.033]$, $p = 0.002$. The diversity ordering again holds: many other mice age
less than one by $-0.010$ $[-0.016, -0.004]$, $p = 0.003$.
The result also does not depend on where the inclusion threshold is drawn, giving $+0.022$ at three
or more sessions per mouse (24 mice), $+0.023$ at four (23 mice), $+0.023$ at five (22 mice) and
$+0.024$ at six (18 mice).

One aspect of the analysis does not transfer between the two cohorts. Repeating the trial-level
split here \emph{lowers} the asymmetry to
$+0.014$ ($p = 0.07$) rather than raising it. Separating the two differences on the same trials:
trial shuffling still inflates the personal decoder by $+0.015$, as on BraiDyn-BC, but scoring one
pooled $\auc$ across sessions instead of averaging per-session values costs $-0.030$
($p < 10^{-4}$) and dominates. The corrected estimate is here the \emph{larger} one, so the
replication is not an artifact of the split.

\begin{table}[t]\centering
\caption{The aging asymmetry on the original cohort and on the external validation dataset. Aging
is $\auc$ lost from early to late test with training held fixed; asymmetry is personal $-$ pooled
aging, positive when the subject's own decoder ages more. Both rows use the same session-held-out
protocol and the same count-matched control. $p$: paired sign-flip permutation.}
\label{tab:allen}
\setlength{\tabcolsep}{3pt}
{\small
\begin{tabular}{lccccc}
\toprule
& Personal & Pooled & Asymmetry (95\% CI) & $p$ & own ages more \\
\midrule
BraiDyn-BC, widefield, 25 mice & & & $+0.017$ $[0.010, 0.024]$ & $5\times10^{-5}$ & 73\% \\
Allen VB, two-photon, 23 mice & $+0.024$ & $+0.001$ & $+0.023$ $[0.010, 0.037]$ & 0.003 & 74\% \\
\bottomrule
\end{tabular}}

\vspace{3pt}
\begin{minipage}{0.95\linewidth}\small\raggedright
\emph{Count-matched control, aging at equal training-event count.} BraiDyn-BC: self $-$ one other
mouse $= +0.010$; many others age less than one on all four events. Allen VB: self $+0.024$, one
$+0.013$, many $+0.003$; self $-$ one $= +0.011$, many $-$ one $= -0.010$.
\end{minipage}
\end{table}

\section{Discussion}
A subject's own decoder overfits idiosyncratic features of its early sessions, some of which later
drift. A decoder that has never seen the subject must instead rely on components of the code that are
shared across individuals and more stable over time. The count-matched control supports this reading
directly: at equal data, other-individual data ages less than the subject's own. Population structure
is therefore not only a resource for cross-subject transfer but also for temporal robustness, and the
early-week accuracy cost of cross-subject decoding is recovered over time as the pooled decoder
remains stable. Unlike within-subject recalibration, this requires no newly labelled data from the
individual at test time -- which is the practical point, since obtaining that data is the expensive
part of deploying a decoder over months.

The effect is small in absolute terms, and we think it is worth being explicit about why we believe
it anyway. It rests on 25 animals and 357 recording sessions, which is a large cohort for a
longitudinal cortex-wide study, and it is supported by five checks that are not variations on one
analysis. A two-block contrast and a day-by-day regression over all fifteen sessions are different
estimators and select the same events. The count-matched control fixes the training-set size, so the
result cannot be a data-volume artifact. Replacing the linear readout with a 1-D CNN or a GRU leaves
it in place, so it is not an artifact of a weak decoder. And holding out whole sessions rather than
trials removes an inflation large enough to eliminate one of the four events entirely -- the effect
reported here is what remains after that correction, not before it. And it reproduced in a second
cohort recorded in another laboratory with a different instrument and a different task.
Taken together, each of these checks could independently have removed the result, and none of them
did.

\begin{figure}[t]
\centering
\fbox{\begin{minipage}{0.95\linewidth}
\textbf{Best practice for longitudinal within- versus cross-subject comparisons}
\begin{enumerate}\itemsep1pt\parskip0pt
\item \textbf{Hold out whole recording sessions, not trials.} Only the within-subject arm can share
a session between train and test, so a trial-level split inflates that arm alone. Here it inflated
within-block accuracy by $+0.013$ (lever pull), $+0.010$ (lick), $+0.009$ (reward) and $+0.000$
$\auc$ (tone), and on lever pull it accounted for the entire apparent effect: $+0.014$
($p = 0.001$) became $+0.003$ ($p = 0.37$).
\item \textbf{Count-match the training sets before comparing arms.} A pooled decoder contains
$\sim$24$\times$ more training events than a personal one. Fix the training-set size to $n$ and fit
\emph{self}, \emph{one} and \emph{many} on the same $n$, scored on the same held-out session.
\item \textbf{Never score an arm on the block it was fit on.} In the count-matched control the
natural implementation scores the self decoder on its own training block -- not a cross-validated
quantity at all, and enough to reverse the ordering the control exists to establish.
\item \textbf{Report each arm's aging separately, not only the difference.} The asymmetry is a
difference of two quantities that can move for different reasons: on BraiDyn-BC it narrows because
the pooled decoder gains, on Allen because the personal decoder loses.
\item \textbf{Average per-session scores rather than pooling one score across sessions.} On the
Allen cohort, scoring a single pooled $\auc$ across sessions instead of averaging per-session values
cost $-0.030$ ($p < 10^{-4}$) and dominated the leakage term it was confounded with.
\end{enumerate}
\end{minipage}}
\label{box:practice}
\end{figure}

\textbf{Limitations.} The clearest limitation is what the four events are. They are not four
independent behaviors but four thresholded channels of one cued lever-pull trial, overlapping to the
point where $100\%$ of reward windows also contain licking, and differing tenfold in event count. We
therefore cannot always say which behavior the effect belongs to, nor can we exclude the idea that
the per-event ranking is a tracker of how much data each decoder had to overfit. A design with
temporally separated, count-matched events would settle both, and is the experiment we would run
next.

A second limitation is spatial scale. We measure at the mesoscale, using parcel-averaged activity,
and while the two-photon replication does reach single neurons it does so only at one depth in one
visual area. Whether the same asymmetry holds for single-neuron codes across cortex is therefore
still open.

A third is temporal and taxonomic scope. The analysis covers a two-week interval in one species, and
nothing here establishes how the asymmetry behaves over the months-to-years horizon on which a
deployed decoder actually has to survive.

Finally, the trial-level inflation of Section~\ref{sec:protocol} is worth stating as a caution
rather than as a limitation of this work, because it is a limitation we do not have. It is large
enough to manufacture an effect that is not there -- on lever pull it manufactured the entire effect
-- and every number we report is measured after holding out whole sessions, not before. We raise it
because any comparison of this shape is exposed to it, and the correction is the one in
Box~\ref{box:practice}.

% Camera-ready only. Uncomment when the paper is de-anonymized -- naming advisors in a
% double-blind submission breaks anonymity.
% \subsubsection*{Acknowledgments}
% We thank Vibha Singhal for detailed feedback on the structure and framing of an earlier draft.

\subsubsection*{Reproducibility}
BraiDyn-BC is publicly available through DANDI:001425. We release the analysis code and a streaming
feature pipeline that reproduces all reported results from the parcellated traces without downloading
the raw imaging data. Section 3 and the released code specify all data splits, permutation seeds, and
hyperparameters. The external validation dataset of Section~\ref{sec:allen} is likewise public, and
the released code streams it from its public bucket and reproduces that section from the same cohort
selection rule.

\bibliography{refs}

\appendix

\section{Extended related work}
\label{app:related}
\textbf{Representational drift and its mitigation.} Drift is characterized within individual
subjects \citep{rule2020,driscoll2017,schoonover2021,rule2019}, and most approaches to mitigating it
also operate within the subject. These include manifold stabilization \citep{degenhart2020},
supervised and self-calibrating recalibration \citep{jarosiewicz2015,wilson2025}, and adversarial or
nonlinear alignment to the subject's own reference day \citep{farshchian2019,karpowicz2025}. Each
method uses later data from the same subject to stabilize the decoder. We instead ask whether data
from other subjects can provide that stability.

\textbf{Cross-subject decoding.} Shared structure across individuals already supports cross-subject
readout through manifold alignment \citep{melbaum2022}, contrastive embeddings \citep{schneider2023},
and, on this cohort, atlas-aligned foundation models \citep{wicat2026}. Multi-participant pooling can
also reduce overfitting: HTNet, for example, pools participant data to avoid fitting a decoder to a
small single-subject sample \citep{peterson2021}, and pooled pretraining across 83 animals followed
by fine-tuning on held-out animals reaches state-of-the-art encoding and decoding on a standardized
cross-subject cohort \citep{neds2025}. Our question differs in two ways. Primarily, prior work
measures accuracy or transfer at a fixed time, while we look at how that accuracy changes over days,
to see whether pooled decoders age more slowly. Further, we keep the training-set size fixed with the
count-matched control, hence being able to delegate the relevant overfitting not to limited data but
to idiosyncratic features that cause drift.

The closest work to ours combines both axes: Meta-AlignNN \citep{metaalign2025} uses the consistency
of neural population dynamics across individuals to build a decoder that remains stable across
subjects, time and tasks, and demonstrates it over two years in three monkeys. That work and ours
point in the same practical direction, and differ in kind rather than in topic. Meta-AlignNN proposes
a method for achieving stability; we measure whether a difference in aging rate exists at all between
a decoder trained on the subject and one trained on others, with the training-set size held fixed so
that the comparison isolates whose data the decoder saw. A measured effect of this kind is what would
justify the design choice such methods already make.

\textbf{Conservation of cortex-wide dynamics.} Cortex-wide activity contains shared,
low-dimensional, movement-related components \citep{musall2019,macdowell2020,macdowell2024}, and
latent dynamics are preserved across animals performing similar behaviour to the point that a decoder
fit to one animal transfers to another \citep{safaie2023}. That conservation is the mechanism a
pooled decoder would exploit, and it is why pooling is possible at all here. It does not, however,
imply that pooling will \emph{resist drift}: a code can be shared across individuals and still move
within each of them at the same rate. Whether the shared component is also the more temporally stable
one is the empirical question we ask.


\section{External validation cohort in full}
\label{app:allen}
Section~\ref{sec:allen} repeats the analysis on the Allen
Visual Behavior two-photon dataset \citep{allenvb2023}. We select containers spanning at least 14
days whose experiments are all in VISp within a single 50\,\textmu m depth bin, which is the
cross-subject correspondence available here; the largest such band, 150--199\,\textmu m, holds 25
mice and 184 experiments. We require at least 15 segmented cells and 40 usable trials per
experiment, leaving 166 experiments and 23 mice with at least four sessions. Sessions are streamed
from S3 rather than downloaded.

As different animals have no neuron correspondence and the longitudinal subset is single-plane,
averaging over the field of view would give a one-dimensional feature. The representation that is
comparable across animals is
therefore a permutation-invariant summary of the population response: per neuron we take the mean
$\Delta F/F$ over $0$--$1$\,s minus the mean over $-0.5$--$0$\,s, then describe the distribution of
those evoked responses by 25 quantiles plus its mean and standard deviation, giving a fixed
27-dimensional feature regardless of how many cells were segmented. The decoded contrast is
\texttt{is\_change}, a real image change against a sham catch trial, which is a stimulus contrast
rather than a reward- or lick-coupled one. Each mouse's sessions are split by date into early and
late halves, and the personal decoder, the pooled decoder, aging, the asymmetry and the
count-matched control are all defined exactly as above.


\section{Day-by-day decay regression}
\label{app:decay}
The block split compares two windows, but the effect should also appear as a continuous trend across
all fifteen days. For each session day we compute the accuracy of the personal decoder and the pooled
decoder and regress each on day; on early days the personal decoder is trained by
leave-one-session-out, so it is never tested on a session it trained on. The personal decoder starts
ahead, leading by about $0.035$ $\auc$ on day 1 as expected for a decoder trained on the same animal.
That lead then erodes: the gap narrows by $0.0014$ $\auc$ per day ($p = 5 \times 10^{-5}$; narrowing
in 65\% of pairs), so by day 15 roughly 60\% of the early-week advantage is gone.

Averaged over the four events, the pooled decoder gains accuracy over days ($+0.0015$ $\auc$/day) as
the sessions become easier -- the mouse is more practiced -- while the personal decoder stays flat
($+0.0000$ $\auc$/day). The personal decoder, fit to idiosyncrasies, is unable to ride that wave, and
the lift it fails to take is what opens the gap. Per event, the personal decoder does lose accuracy
for lick ($-0.0010$ $\auc$/day) and reward ($-0.0016$ $\auc$/day), and these are the two events where
the narrowing is individually significant ($p = 0.007$ and $p = 0.001$; Fig.~\ref{fig:decay}).

\begin{figure}[t]\centering
\includegraphics[width=0.98\linewidth]{figs/fig_decay.pdf}
\caption{Decoding accuracy across the 15-day protocol for the personal decoder (own early data) and
the pooled decoder (other mice's early data), with per-day points and linear fits. The personal
decoder starts ahead, and the pooled decoder gains more over days, so the gap between them narrows.
The narrowing is significant for lick and reward, the two events where the personal decoder also
loses accuracy over days, and is not significant for lever pull or tone. The per-event gap slope and
its $p$ (paired sign-flip over mice) are shown.}
\label{fig:decay}\end{figure}

\end{document}
